\documentclass[11pt]{article}

\usepackage[margin=1in]{geometry}
\usepackage{amsmath,amssymb,amsthm,mathtools}
\usepackage{enumitem}
\usepackage{hyperref}

\setlength{\parskip}{0.5em}
\setlength{\parindent}{0pt}

\newcommand{\R}{\mathbb{R}}
\newcommand{\K}{\mathcal{K}}
\newcommand{\eps}{\varepsilon}
\newcommand{\bstar}{f^\ast}
\newcommand{\xstar}{x^\ast}

\begin{document}

\begin{center}
{\Large \textbf{DSC 243 --- Homework 1}} \\
\vspace{0.5em}
{\large Convex Quadratics, Chebyshev Acceleration, and Krylov Methods}
\end{center}

\textbf{Instructions.}
Use the notation from the notes. Unless explicitly stated otherwise, assume
\begin{align*}
&f(x) = \frac{1}{2} x^\top A x - b^\top x,\\
&A = A^\top \succeq 0,\\
&b \in \mathrm{range}(A),\\
&\alpha = \lambda_{\min}(A),\\
&\beta = \lambda_{\max}(A),\\
&\text{and when } \alpha > 0,\,\, \kappa = \beta/\alpha.
\end{align*}
You may quote results proved in class or in the notes, but every nontrivial step must be justified.

\section*{Part I: Theory}

\begin{enumerate}[label=\textbf{T\arabic*.}, leftmargin=0pt, itemindent=2.5em]

\item \textbf{Support-dependent finite termination.}
Let $A\succeq 0$, and let $\lambda_1,\dots,\lambda_m$ be the distinct nonzero eigenvalues of $A$ with corresponding orthonormal eigenvectors $v_1,\dots,v_m$. Suppose the initial residual satisfies
\[
r_0 \in \operatorname{span}\{v_{i_1},\dots,v_{i_r}\}
\]
for some subset of $r$ distinct eigendirections. Show that the Krylov method terminates in at most $r$ iterations.

\noindent
[\textbf{Hint:} Use that every point in $x_0+\K_k(A, r_0)$ can be written as $x_0 + q(A) r_0$ for some polynomial $q$ of degree at most $k-1$.]


\item \textbf{Repeated Chebyshev cycles.}
In this problem, assume $A$ is positive definite and $\sigma(A)\subset[\alpha,\beta]\subset(0,\infty)$.
Fix an integer $s\ge 1$, and let $\eta_1,\dots,\eta_s$ be the degree-$s$ Chebyshev stepsizes from Theorem 3.1. Define the polynomial
\[
p_s(\lambda)=\prod_{j=1}^s(1-\eta_j\lambda).
\]
Repeat this same $s$-step schedule for $M$ epochs, so the total number of iterations is $N=Ms$.
\begin{enumerate}[label=(\alph*)]
\item Show that the error after $N$ iterations is
\[
e_N=p_s(A)^M e_0.
\]
\item Prove that
\[
f(x_N)-\bstar \le \left(\max_{\lambda\in[\alpha,\beta]}|p_s(\lambda)|^2\right)^M \bigl(f(x_0)-\bstar\bigr).
\]
\item Deduce the bound
\[
f(x_N)-\bstar
\le
\frac{1}{T_s(\sigma)^{2M}}\bigl(f(x_0)-\bstar\bigr)
\le
4^M\left(\frac{\sqrt{\kappa}-1}{\sqrt{\kappa}+1}\right)^{2N}\bigl(f(x_0)-\bstar\bigr).
\]
\item Using the exact bound
\[
f(x_N)-\bstar \le \frac{1}{T_s(\sigma)^{2M}}\bigl(f(x_0)-\bstar\bigr),
\]
derive an iteration-complexity guarantee for the restarted scheme: how large must $N=Ms$ be in order to ensure
\[
f(x_N)-\bstar \le \eps?
\]
Express your answer in terms of $\kappa$, $s$, $\eps$, and $f(x_0)-\bstar$.
\item The one-shot degree-$N$ Chebyshev schedule satisfies
\[
f(x_N)-\bstar \le \frac{1}{T_N(\sigma)^2}\bigl(f(x_0)-\bstar\bigr).
\]
Derive the corresponding iteration-complexity estimate for this one-shot scheme.
\item Compare the two complexity estimates. In particular, identify explicitly where the restarted method loses efficiency relative to the one-shot degree-$N$ schedule.
\end{enumerate}

\item \textbf{(Optional) Krylov convergence under clustered eigenvalues.}
In this problem, assume $A$ is positive definite and that its eigenvalues lie in a union of $r$ disjoint intervals
\[
[\ell_1,u_1],\ [\ell_2,u_2],\ \dots,\ [\ell_r,u_r]\subset[\alpha,\beta].
\]
Define the cluster centers and half-widths by
\[
c_j:=\frac{\ell_j+u_j}{2},
\qquad
\delta_j:=\frac{u_j-\ell_j}{2},
\qquad j=1,\dots,r.
\]
Also define
\[
C_{\mathrm{clust}}
:=
\max_{1\le j\le r}
\left(
\frac{\delta_j}{c_j}
\prod_{i\ne j}\left(1+\frac{u_j}{c_i}\right)
\right).
\]
\begin{enumerate}[label=(\alph*)]
\item Consider the degree-$r$ polynomial
\[
\pi_r(\lambda):=\prod_{j=1}^r\left(1-\frac{\lambda}{c_j}\right).
\]
Show that $\pi_r(0)=1$ and that, for every $\lambda\in[\ell_j,u_j]$,
\[
\left|\pi_r(\lambda)\right|
\le
\frac{\delta_j}{c_j}
\prod_{i\ne j}\left(1+\frac{u_j}{c_i}\right).
\]

\item Deduce that
\[
\max_{\lambda\in \sigma(A)} |\pi_r(\lambda)| \le C_{\mathrm{clust}}.
\]

\item Let $q_t$ be any degree-$t$ polynomial with $q_t(0)=1$, and define
\[
p_{r+t}(\lambda):=\pi_r(\lambda)\,q_t(\lambda).
\]
Explain why $p_{r+t}(0)=1$ and why the Krylov method after $r+t$ iterations satisfies
\[
f(x_{r+t})-\bstar
\le
\left(\max_{\lambda\in\sigma(A)} |\pi_r(\lambda)|^2\right)
\left(\max_{\lambda\in[\alpha,\beta]} |q_t(\lambda)|^2\right)
\bigl(f(x_0)-\bstar\bigr).
\]

\item Now choose $q_t$ to be the degree-$t$ shifted Chebyshev polynomial from Section 3. Deduce the bound
\[
f(x_{r+t})-\bstar
\le
\frac{C_{\mathrm{clust}}^2}{T_t(\sigma)^2}\bigl(f(x_0)-\bstar\bigr)
\le
4C_{\mathrm{clust}}^2
\left(\frac{\sqrt{\kappa}-1}{\sqrt{\kappa}+1}\right)^{2t}(f(x_0)-\bstar).
\]

\item Interpret the result. What happens when the clusters collapse to points, i.e. $\delta_j\to 0$ with the centers $c_j$ fixed? How does this relate to the exact finite-termination theorem for $r$ distinct eigenvalues?
\end{enumerate}

\noindent
[\textbf{Hints:}
For part (a), if $\lambda\in[\ell_j,u_j]$, then
\[
\left|1-\frac{\lambda}{c_j}\right|=\frac{|\lambda-c_j|}{c_j}\le \frac{\delta_j}{c_j}.
\]
For $i\ne j$, use
\[
\left|1-\frac{\lambda}{c_i}\right|\le 1+\frac{\lambda}{c_i}\le 1+\frac{u_j}{c_i}.
\]
For part (d), combine parts (b) and (c) with the standard Chebyshev bound on $[\alpha,\beta]$.]



\end{enumerate}

\section*{Part II: Experiments}

\begin{enumerate}[label=\textbf{E\arabic*.}, leftmargin=0pt, itemindent=2.5em]

\item \textbf{CG termination and spectral support.}
Take $d=120$ and let
\[
A=\operatorname{diag}(\lambda_1,\dots,\lambda_{20},0,\dots,0),
\]
where $\lambda_1,\dots,\lambda_{20}$ are $20$ distinct positive numbers of your choice. Let
\[
x^\ast=\mathbf{1}\in\R^d,\qquad b=Ax^\ast,
\]
so that $x^\ast$ is a minimizer.
\begin{enumerate}[label=(\alph*)]
\item Choose $x_0$ so that the initial residual $r_0=b-Ax_0=A(x^\ast-x_0)$ is supported on exactly $r=3$ eigendirections corresponding to nonzero eigenvalues, and run CG. Plot $\|r_k\|/\|r_0\|$.
\item Repeat with $r_0$ supported on exactly $r=6$ nonzero eigendirections.
\item Repeat with a generic initial point $x_0$ such that the resulting residual $r_0=b-Ax_0$ has nonzero components along all $20$ nonzero eigendirections.
\item Compare the observed termination behavior with the prediction from Problem T1.
\end{enumerate}

\item \textbf{One-shot Chebyshev versus restarted Chebyshev.}
Take $d=200$ and
\[
A=\operatorname{diag}\left(1,\,1+\frac{99}{d-1},\,1+2\frac{99}{d-1},\,\dots,\,100\right),
\]
so that $A$ is positive definite with $\kappa=100$. Let
\[
x^\ast=\mathbf{1}\in\R^d,\qquad b=Ax^\ast.
\]
Let $N=120$, and use the same generic random initial point $x_0$ for every method.
\begin{enumerate}[label=(\alph*)]
\item Run one-shot degree-$120$ Chebyshev acceleration.
\item For each cycle length $s\in\{5,10,20,30\}$, run restarted Chebyshev with repeated $s$-step cycles until $N$ iterations are used.
\item Plot the relative suboptimality
\[
\frac{f(x_k)-f^\ast}{f(x_0)-f^\ast}
\]
of all methods on the same graph.
\item Report which restart length performs best empirically, and compare this with the theory from Problem T2.
\end{enumerate}

\item \textbf{CG under clustered spectra.}
Fix $d=180$ and cluster centers
\[
c_1=2,\qquad c_2=20,\qquad c_3=80.
\]
For each $\gamma\in\{10^{-1},\,10^{-2},\,10^{-3}\}$, construct a diagonal positive definite matrix whose eigenvalues consist of $60$ points in each interval
\[
\bigl[c_j-\gamma c_j,\ c_j+\gamma c_j\bigr],
\qquad j=1,2,3.
\]
For example, you may take the eigenvalues to be uniformly spaced within each interval. For each such matrix, set
\[
x^\ast=\mathbf{1}\in\R^d,\qquad b=Ax^\ast,
\]
and use the same generic random initial point $x_0$ for all runs.
\begin{enumerate}[label=(\alph*)]
\item For each choice of $\gamma$, run CG and plot
\[
\frac{f(x_k)-f^\ast}{f(x_0)-f^\ast}
\]
versus $k$ on a semilog scale.
\item On the same plots, overlay the standard worst-case Chebyshev/Krylov bound based only on the condition number $\kappa$ of that matrix.
\item Also overlay the clustered-spectrum bound predicted by Problem T3: for $k\ge 3$, use the constant $C_{\mathrm{clust}}$ from T3 and the estimate
\[
f(x_k)-f^\ast
\le
4C_{\mathrm{clust}}^2
\left(\frac{\sqrt{\kappa}-1}{\sqrt{\kappa}+1}\right)^{2(k-3)}
\bigl(f(x_0)-f^\ast\bigr).
\]
\item Compare the three curves. Does the empirical CG behavior track the clustered-spectrum prediction more closely as the clusters become tighter?
\item Briefly explain how the experiment illustrates the transition from generic $O(\sqrt{\kappa}\log(1/\eps))$ behavior toward near finite termination.
\end{enumerate}

\end{enumerate}

\end{document}
